\chapter{Exit codes}

\section{Check the exit code of the previous command}

you can check the exit code of previous commands using the dollar sign. If it's 0, it means
success. Otherwise, it means failure. The number might give you a clue on the failure actually 
was about.

\lstinputlisting[
    language=bash,
    caption={5-bash.sh},
    captionpos=t,
    label={lst:5-bash},
    breaklines=true]
{../5-bash/5-bash.sh}

\section{Redirects}

Redirects allow us to save the output of a command directly into a file.

\lstinputlisting[
    language=bash,
    caption={5-bash.sh},
    captionpos=t,
    label={lst:5.1-bash},
    breaklines=true]
{../5-bash/5.1-bash.sh}

Something to consider is that since the file doesn't follow a strict path, the script will create
the file in your current path where you executed the script. 

\section{Careful with \$?}

Something you have to remember is that \$? will only give you the exit code of the last command
so if you want to check the first command, you have to store the value in another variable

\lstinputlisting[
    language=bash,
    caption={5-bash.sh},
    captionpos=t,
    label={lst:5.2-bash},
    breaklines=true]
{../5-bash/5.2-bash.sh}

\section{Manually exit}

If you didn't want to store the exit value in a variable, which can get convoluted real quick,
you can manually exit your script with a custom code.

\lstinputlisting[
    language=bash,
    caption={5-bash.sh},
    captionpos=t,
    label={lst:5.3-bash},
    breaklines=true]
{../5-bash/5.3-bash.sh}

\section{The final version of the 5-bash.sh script}

\lstinputlisting[
    language=bash,
    caption={5-bash.sh},
    captionpos=t,
    label={lst:5.4-bash},
    breaklines=true]
{../5-bash/5.4-bash.sh}
